% Table 12: Limitations summary
\begin{table}[htbp]
\centering
\caption{Limitations of the HiddenCourtesyMerge-Sim benchmark.}
\label{tab:limitations}
\small
\begin{tabular}{p{3.4cm}p{9.0cm}}
\toprule
\textbf{Limitation} & \textbf{Detail} \\
\midrule
Simulation-only scope
  & All results are from highway-env \texttt{merge-v0}; no real-world or
    human-driver validation is performed or claimed. \\
Synthetic courtesy
  & Courtesy types are operationalised via IDM/MOBIL parameter ranges; they
    do not model empirical human psychology or empirical driving data. \\
Binary courtesy only
  & A two-class latent variable is a controlled simplification; real driver
    behaviour is continuous and multi-dimensional. \\
Heuristic action table
  & The evaluated policies use a hand-designed V-shaped caution heuristic,
    not a value-optimal POMDP policy. The oracle upper bound is for this
    table only. \\
Two-action restriction
  & The heuristic policies restrict actions to \texttt{IDLE} and
    \texttt{SLOWER}; \texttt{FASTER} and lateral actions remain in the
    environment action space but are not used by the evaluated non-random
    policies. \\
Reward penalty weight
  & The close-call coefficient 0.25 was chosen by manual calibration to
    penalise near-conflict timesteps. Replay-based sensitivity over
    $w\in\{0.10,0.25,0.50,1.00\}$ preserves the policy reward ordering
    (Table~\ref{tab:penalty_ablation}), but closed-loop behaviour under
    alternative rewards was not rerun. \\
Static courtesy
  & Courtesy is sampled once per episode and held fixed through IDM/MOBIL
    parameters. This is appropriate for the controlled benchmark but does not
    model drivers whose style changes mid-encounter. \\
QMDP excluded
  & A QMDP approximation was implemented (16-state tabular Q-learning) but
    excluded from the primary evaluation due to insufficient Q-table coverage
    ($\ge$50\% of states unvisited with $n{=}50$ training episodes). \\
Belief vs.\ rule underpowered
  & The belief policy does not significantly outperform the rule baseline
    at $n{=}300$; $\approx$1{,}960 episodes would be needed for 80\% power. \\
Calibration imbalance
  & The observation model was calibrated on 290 cooperative and 1{,}519
    non-cooperative in-window steps (5.2:1 ratio). Cooperative Gaussian
    estimates carry higher uncertainty. \\
Non-interaction episodes
  & 778 of 3{,}000 episodes (26\%) contain no merge interaction window;
    the merge vehicle never entered the 2--40\,m zone and no belief update
    occurred in those episodes. \\
Weak discriminators
  & Urgency (1/TTC) has a separation score of 0.14 (weak) and was dropped
    from the final model; absolute relative distance is borderline (0.48).
    The final filter relies primarily on merge vehicle speed (\emph{mvs},
    score 1.32) and merge vehicle acceleration (\emph{mva}, score 0.81). \\
Sequential episode split
  & Train/validation/test splits are assigned by episode index (70/15/15)
    rather than by stratified shuffle.
    If the data-generating process drifts over the 3{,}000-episode run
    (e.g.\ due to highway-env RNG non-stationarity), later splits may have
    a systematically different state distribution than earlier ones. \\
Data-generating policy
  & The 3{,}000-episode dataset was generated with the \texttt{belief\_policy}
    as ego driver.
    The calibration may underrepresent state regions avoided by
    \texttt{belief\_policy}; robustness to alternative generating policies
    (e.g.\ constant-speed or rule-based ego) has not been verified and remains
    an open question. \\
Belief regularisation
  & The canonical $\lambda{=}0.08$ setting is close to the best final-step
    accuracy but is not optimal for NLL or ECE in Table~\ref{tab:lambda_ablation}.
    Sensitivity of closed-loop policy outcomes to changing $\lambda$ remains
    untested. \\
Short horizon
  & Episodes are capped at 60 steps ($\Delta t{=}0.2$\,s each, i.e., 12\,s).
    The brief interaction window limits the number of belief updates available
    before the episode ends. \\
\bottomrule
\end{tabular}
\end{table}
